% chapters/14/section-06.tex
\section{Aset dan Hak Pakainya}
\label{sec:hak-aset}

\subsection{Jenis Lisensi yang Perlu Dikenal}

Aset yang ditemukan di internet tidak selalu boleh dipakai. Sebelum memakai
aset apa pun, periksa lisensinya. Tiga jenis lisensi yang paling sering
dijumpai pada aset permainan diperlihatkan pada
Tabel~\ref{tab:lisensi-aset}.

\begin{table}[htbp]
\centering
\caption{Tiga jenis lisensi aset yang paling sering dijumpai}
\label{tab:lisensi-aset}
{\setlength{\tabcolsep}{4pt}\footnotesize%
\begin{tabularx}{\textwidth}{@{}L{3.0cm} L{4.4cm} Y@{}}
\toprule
\headrow \textbf{Lisensi} & \textbf{Boleh dipakai untuk} & \textbf{Kewajiban} \\
\midrule
\istilah{Public domain} atau \istilah{CC0} &
Segala keperluan, termasuk yang komersial. &
Tidak ada. Cara yang paling bebas, dan paling aman untuk proyek mata kuliah. \\
\addlinespace
\istilah{Creative Commons} dengan atribusi (CC~BY) &
Segala keperluan, termasuk yang komersial. &
Wajib menyebutkan pembuatnya dan lisensinya. Untuk keperluan mata kuliah ini,
cukup dicantumkan pada laporan. \\
\addlinespace
\istilah{Creative Commons} tanpa penggunaan komersial (CC~BY-NC) &
Keperluan pendidikan, termasuk tugas mata kuliah. &
Wajib menyebutkan pembuatnya, dan tidak boleh dipakai untuk produk yang
dijual. \\
\bottomrule
\end{tabularx}}
\tabsource{Bacalah lisensi setiap aset pada halaman asalnya. Sebagian aset
menyertakan berkas lisensi tersendiri di dalam paketnya.}
\end{table}

\begin{peringatanbox}
Aset yang ditemukan melalui mesin pencari gambar belum tentu boleh dipakai.
Hak cipta berlaku sejak karya itu dibuat, walaupun tidak ada keterangan apa
pun pada berkasnya. Kebiasaan yang aman: pakai hanya aset dari tempat yang
menyatakan lisensinya dengan jelas, dan catat sumbernya ketika mengimpor.
\end{peringatanbox}

\subsection{Mencatat Sumber Aset}

Catatan sumber aset sebaiknya dibuat sejak awal, dalam bentuk yang sederhana.
Empat hal perlu dicatat untuk setiap aset: nama berkasnya, sumbernya, jenis
lisensinya, dan keperluan atribusinya.

Catatan itu berguna untuk tiga hal: memenuhi kewajiban atribusi, memudahkan
Anda mencari penggantinya apabila diperlukan, dan menjadi bukti kejujuran
akademik pada laporan proyek akhir Anda. Dua contoh laporan akademik yang
memakai Unity dan mencantumkan asetnya secara terbuka dapat dilihat pada
\cite{uny-katelu} dan \cite{undip-media}.

\begin{tipsbox}
Kebiasaan yang sama berlaku untuk bunyi, huruf, dan potongan kode yang Anda
ambil dari internet. Buku ini sendiri memakai kebiasaan itu: setiap gambar yang
bukan karya penulis diberi keterangan sumbernya di bawah gambarnya, misalnya
pada Gambar~\ref{fig:editor-unity} dan Gambar~\ref{fig:editor-fisika}.
\end{tipsbox}

\subsection{Ringkasan Perkembangan}

\begin{table}[htbp]
\centering
\caption{Bagian yang ditambahkan pada aset visual pada bab berikutnya}
\label{tab:perkembangan-aset}
{\setlength{\tabcolsep}{4pt}\footnotesize%
\begin{tabularx}{\textwidth}{@{}C{1.6cm} L{5.0cm} Y@{}}
\toprule
\headrow \textbf{Bab} & \textbf{Yang ditambahkan} & \textbf{Akibatnya} \\
\midrule
Bab~XV & Bunyi dan umpan balik interaksi &
Peristiwa permainan---lompatan, pengambilan benda, kalah---memperoleh tanggapan
bunyi. \\
\addlinespace
Bab~XVI & Pengujian prototipe dan presentasi &
Aset yang telah dipasang dinilai mutunya oleh pemain lain, dan hasilnya
menjadi dasar perbaikan. \\
\bottomrule
\end{tabularx}}
\end{table}

Dengan gambaran itu, tampak bahwa bab ini melengkapi bagian pertama \CPMK{4},
yaitu penyajian visual. Yang tersisa adalah penyajian bunyi, dan pengujian atas
keduanya secara menyeluruh.
